(1) 约束 \(z=ar^2\) 导致 \(\dot z=2ar\dot r\)，且 \(\dot\phi=\Omega\)，故 \(L=\frac m2[(1+4a^2r^2)\dot r^2+\Omega^2r^2]-mg ar^2\)。令 \(A(r)=1+4a^2r^2\)，Euler--Lagrange 方程为 \(mA\ddot r+4ma^2r\dot r^2+m(2ga-\Omega^2)r=0\)。其中速度平方项来自径向动能的坐标依赖；忽略它仅在线性化附近成立。

(2) 以径向约化坐标写 \(L=\frac12mA(r)\dot r^2-U_{\rm eff}(r)\)，其中 \(U_{\rm eff}=m(ga-\Omega^2/2)r^2\)。定义 \(\Omega_c=\sqrt{2ga}\)，若 \(\Omega<\Omega_c\)，轴心有效势严格正曲率，径向受约束稳定；若 \(\Omega>\Omega_c\)，轴心线性不稳定；恰在临界时平面型有效势全零，任一常半径均满足径向方程，不能称其为通常意义的非线性回复稳定。

(3) 在轴心 \(A(0)=1\)，线性方程是 \(\ddot r+(2ga-\Omega^2)r=0\)，故 \(\omega_r=\sqrt{2ga-\Omega^2}\)（仅亚临界为实）。非零半径静平衡要求 \((2ga-\Omega^2)r_e=0\)。因此除恰临界整条平坦平衡族外，不存在通常可稳定的有限 \(r_e>0\) 分岔支。若有人根据不同的旋转圆环模型写出非零有限半径稳定解，就违反了本抛物面 \(z=ar^2\) 的具体几何约束。

(4) 惯性系中珠子的机械能为 \(E_{\rm in}=\frac m2[A(r)\dot r^2+\Omega^2r^2]+mg ar^2\)；外部旋转机构为维持 \(\dot\phi=\Omega\) 持续施加方位约束力矩，其功率为 \(\Omega\,d(m\Omega r^2)/dt=2m\Omega^2r\dot r\)，恰与径向移动引起的惯性系能量变化相容。Jacobi 型常量 \(E_{\rm in}-\Omega L_z=\frac12mA\dot r^2+U_{\rm eff}\) 守恒。当 \(\Omega>\Omega_c\) 时 \(U_{\rm eff}\to-\infty\)，珠可向无限半径滚落；此时理想无限大刚性曲面与固定转速驱动假设的物理代价不能忽略。
